% Chapitre 8 — Agent et orchestration
\chapter{Agent conversationnel et orchestration LangGraph}
\label{ch:agent}

\section{Motivation}

Les endpoints REST individuels (\api{/explain/by-cin}, etc.) requièrent une connaissance technique de l'API. Les analystes crédit préfèrent une interaction en \textbf{langage naturel} :
\begin{quote}
\textit{« Analyse le CIN 88710263 en séquentiel »} \\
\textit{« Compare les modèles pour le client 95822412 »} \\
\textit{« Fais un rapport complet pour le CIN 88710263 »}
\end{quote}

L'agent conversationnel traduit ces requêtes en appels orchestrés aux modèles ML et au module RAG.

\section{Choix technologiques}

\begin{table}[H]
  \centering
  \caption{Comparaison des frameworks d'orchestration}
  \begin{tabular}{|l|p{4cm}|p{4cm}|p{4cm}|}
    \hline
    & \textbf{LangChain Agents} & \textbf{LangGraph} & \textbf{LangFlow} \\
    \hline
    Contrôle du flux & Faible & \textbf{Élevé (graphe)} & Visuel \\
    \hline
    Debug / traçabilité & Moyen & \textbf{États explicites} & Moyen \\
    \hline
    Production & Variable & \textbf{Recommandé} & Prototypage \\
    \hline
    Choix projet & Écarté seul & \textbf{Retenu} & Non requis \\
    \hline
  \end{tabular}
\end{table}

\textbf{LangGraph} a été retenu pour son modèle de \textbf{graphe d'états typé} (\texttt{AgentState}), facilitant le respect du principe SRP (un nœud = une responsabilité).

\section{État partagé — AgentState}

\begin{lstlisting}[caption={Structure AgentState (extrait)},label={lst:agentstate}]
class AgentState(TypedDict, total=False):
    session_id: str
    user_message: str
    cin: str | None
    credit_id: int | None
    intent: Intent
    route: Literal["classic", "sequential", "graph", ...]
    classic_result: dict | None
    sequential_result: dict | None
    graph_result: dict | None
    rag_sources: list[dict]
    report_markdown: str
    final_answer: str
    model_selected: str | None
    errors: list[str]
\end{lstlisting}

\section{Détection d'intent}

Heuristique par mots-clés (extensible vers un classificateur LLM) :

\begin{table}[H]
  \centering
  \caption{Règles de détection d'intent}
  \begin{tabular}{|l|p{8cm}|}
    \hline
    \textbf{Intent} & \textbf{Mots-clés déclencheurs} \\
    \hline
    compare\_models & comparer, compare, vs, versus \\
    \hline
    full\_report & rapport, report, complet, full \\
    \hline
    graph\_score & graphe, graph, gnn, graphsage, réseau \\
    \hline
    sequential\_score & séquentiel, lstm, gru, transaction \\
    \hline
    classic\_score & (défaut) \\
    \hline
  \end{tabular}
\end{table}

\section{Extraction CIN}

Expression régulière : \texttt{\textbackslash b(\textbackslash d\{6,32\})\textbackslash b}

Exemples valides : \texttt{88710263}, \texttt{95822412}.

\section{Endpoint /chat}

\begin{lstlisting}[caption={Exemple requête /chat}]
POST /chat
{
  "session_id": "demo-1",
  "message": "Analyse le CIN 88710263 en séquentiel"
}
\end{lstlisting}

Réponse :
\begin{lstlisting}[caption={Exemple réponse /chat}]
{
  "session_id": "demo-1",
  "model_selected": "sequential",
  "cin": "88710263",
  "intent": "sequential_score",
  "answer": "Intent détecté: sequential_score | CIN: 88710263\n\n## Résumé\n..."
}
\end{lstlisting}

\section{Gestion des erreurs}

\begin{itemize}[leftmargin=*]
  \item CIN absent $\rightarrow$ message guidant l'utilisateur ;
  \item Échec modèle $\rightarrow$ erreur capturée dans \texttt{errors[]}, rapport fallback ;
  \item Ollama indisponible $\rightarrow$ rapport Markdown déterministe sans LLM ;
  \item \texttt{DISABLE\_LLM=1} lors des appels internes orchestrateur (performance).
\end{itemize}

\section{Mémoire conversationnelle}

Une mémoire session simple (\texttt{\_chat\_history}) stocke les échanges pour continuité UI. Une évolution vers Redis ou PostgreSQL est prévue pour la production multi-utilisateurs.

\section{LangChain vs LangGraph — rôle respectif}

\begin{itemize}[leftmargin=*]
  \item \textbf{LangChain} : wrappers LLM (\texttt{ChatOllama}), messages System/Human ;
  \item \textbf{LangGraph} : orchestration déterministe du workflow, indépendante du choix LLM.
\end{itemize}

Cette séparation permet de remplacer Ollama par un LLM cloud (Azure OpenAI, etc.) sans modifier le graphe.

\section{Tests de l'agent}

\begin{table}[H]
  \centering
  \caption{Scénarios de test agent}
  \begin{tabular}{|p{5cm}|p{4cm}|p{4cm}|}
    \hline
    \textbf{Message} & \textbf{Intent attendu} & \textbf{Modèle} \\
    \hline
    Analyse CIN X & classic\_score & classic \\
    \hline
    CIN X en séquentiel & sequential\_score & sequential \\
    \hline
    CIN X graphe & graph\_score & graph \\
    \hline
    Rapport complet CIN X & full\_report & les 3 \\
    \hline
    Compare modèles CIN X & compare\_models & les 3 \\
    \hline
  \end{tabular}
\end{table}
